\documentclass[11pt,letterpaper]{article}

\usepackage[utf8]{inputenc}
\usepackage[T1]{fontenc}
\usepackage{lmodern}
\usepackage{geometry}
\usepackage{microtype}
\usepackage{amsmath,amssymb,amsthm}
\usepackage{booktabs}
\usepackage{array}
\usepackage{enumitem}
\usepackage{xcolor}
\usepackage[hidelinks]{hyperref}

\geometry{margin=1.15in}
\setlist{itemsep=2pt, topsep=4pt}
\renewcommand{\arraystretch}{1.25}

\theoremstyle{plain}
\newtheorem{theorem}{Theorem}
\newtheorem{proposition}[theorem]{Proposition}
\theoremstyle{definition}
\newtheorem{definition}[theorem]{Definition}
\newtheorem{posit}[theorem]{Posit}
\newtheorem{open}[theorem]{Open}

\newcommand{\Rb}{\mathbb{R}}
\newcommand{\Eb}{\mathbb{E}}

\title{\vspace{-1.2cm}\textbf{Observer-Limited Geometric Cognition}\\[0.5em]
\large The identification problem, and what has to hold for the theory\\
to say anything a flexible embedding does not}

\author{Andrew H. Bond\\
\small Department of Computer Engineering, San Jos\'e State University\\
\small \texttt{andrew.bond@sjsu.edu}}

\date{\small Draft 0.1, September 2026\\[0.4em]
\small Every statement is labeled \emph{proved}, \emph{defined}, \emph{posited},
or \emph{open}.}

\begin{document}
\maketitle

\begin{abstract}
\noindent
Observer-Limited Geometric Cognition holds that an agent represents a task under
finite resolution, evaluates represented consequences in a geometry of its own,
and selects behavior by bounded search. Preference is then a derived object
rather than a primitive. The formal core of this picture is not new. Geometric
Evaluation Theory supplies the evaluation object and proves that a rank budget
coarsens distinctions, that the induced preference is a semiorder with threshold
equal to the budget, and that a budget change reverses preference between fixed
actions. Observation Theory supplies the prior layer, in which what a consumer
can read is a property of the pair and not of the system. This paper argues that
the open problem is neither geometry nor dynamics but identification. A
representation, a metric, a reference, an admissible set and a search policy
together form a family expressive enough to fit any finite choice record, so a
theory that fits every record predicts nothing. We state that as a proposition
rather than as a worry, and give the conditions under which the account
carries content. The conditions are modular. Each of the observer's four
components must be identified on evidence disjoint from the choices the theory
is asked to predict. A cold reread of the first phase's registration
then shows that this condition is not sufficient. Where the class a pair is
assigned to is itself read off the fitted model, the graded contrast is an
algebraic identity in two agreement rates and measures nothing else, verified
here against six cells to within the instrument's own reproducibility, so
content is not monotone in identification accuracy and we can bound neither end.
A rank budget carries a further condition that modular identification does not
imply, since the retained subspace is identified only where the workload moment
has a spectral gap, and a design that places the ideal at the centre of its own
consequence distribution has no gap and recovers a direction fixed by sampling
noise. What survives is that identification is the step
that fails in practice, since two of three candidate evaluators could not be
calibrated at all. Identifying a budget is also not the same as changing one,
and we report three attempts to change one in a language-model evaluator, none
of which produced a resolution effect. A projected stimulus is not a reduced
budget, removing the evaluator's deliberation removes the instrument rather than
narrowing it, and a load the instrument does survive moved nothing we could
measure. None of the three included a check that the intervention reduced
effective rank, which is why we report them as an open question rather than as
evidence against the prediction. We restate the five
predictions of the programme as
consequences of those conditions, restate the seven falsifiers as conditions on
the record rather than on the author's judgment, and give a phased programme in
which each phase retires exactly one degree of freedom. The theory's own
epistemic status is that the evaluation layer is proved and measured on
artificial evaluators, the observer layer is measured in other domains, and the
cognitive synthesis is posited.
\end{abstract}

\vspace{0.6em}
\noindent\textbf{Keywords:} geometric cognition, observer-relative
representation, resolution budget, preference reversal, identification,
bounded search

\section{What this paper claims and does not claim}

The thesis behind this paper is that intelligent agents act by searching over
observer-relative, resource-limited geometries of possible states. This paper
does not argue for that thesis. It asks the prior question of what would have to
be true for the thesis to be testable at all, and answers it.

Three things are already established elsewhere and are used here without
reproof. Geometric Evaluation Theory (GET) proves that unbudgeted distinctions
form an equivalence, that a rank budget coarsens it, that the induced preference
is a weak order at zero threshold and a semiorder at positive threshold with the
threshold fixed by the budget, and that a rank budget change reverses preference
between fixed actions. Observation Theory (OT) establishes in several domains
that what a consumer can recover is a property of the pair formed by the system
and the consumer rather than of the system alone. A domain exposition of
geometric cognition exists and reports measurements on artificial agents.

One thing is not established, and this paper is about it. The domain exposition
measured cognition through a representation that it also chose. Its admission
filter recorded the resulting circularity and refused the instrument, on the
ground that no structure label independent of the embedding was available. That
refusal is correct and it is the starting point here. A geometric account of
cognition whose representation is fitted to the same choices it explains is not
a weak theory. It is not yet a theory.

\section{The evaluation object and the observer above it}

\begin{definition}[Evaluation object, after GET]
An evaluation object is a tuple
\[
E = (X, A, C, G, B, K)
\]
with states $X$ carrying a belief measure $\mu$, actions $A$, a consequence map
$C \colon A \times X \to Y$, a pointed geometry $G$ consisting of a Riemannian
metric on $Y$ and an ideal $t$, a resolution budget $B$, and an admissible set
$K$ that may depend on the state or on the menu.
\end{definition}

The commitment that makes the object non-trivial is the split. The world
supplies $(X, A, C)$ and the evaluator supplies $(G, B, K)$. Without it an
evaluator free to choose $C$ after seeing the desired ordering represents any
finite order by construction, and representability is a tautology. With it,
representation and identification become questions with answers.

Evaluation is by expected squared geodesic distance from the ideal, and cost is
that distance squared,
\[
d(a) = \bigl[\, \Eb_\mu\, d_G\bigl(C(a,X), t\bigr)^2 \,\bigr]^{1/2},
\qquad c(a) = d(a)^2 .
\]
A metric alone orders nothing. Ordering appears only relative to a reference,
which separates the question of which differences the agent is sensitive to,
carried by $G$, from the question of where those differences are measured from,
carried by $t$.

Under a rank budget $k$ the evaluator retains the top-$k$ eigenspace of the
workload-weighted second moment,
\[
M_\omega = \Eb_\omega\bigl[\, G^{1/2}(C-t)(C-t)^{\!\top} G^{1/2} \,\bigr],
\qquad
d_k(a) = \bigl[\, \Eb_\mu \,\|\Pi_k G^{1/2}(C(a,X)-t)\|^2 \,\bigr]^{1/2}.
\]
Two consequences that differ only in discarded directions are then
indistinguishable to the evaluator. This is the structural claim that separates
the account from a noise model. A budget does not blur a complete
representation. It removes directions.

\begin{definition}[Observer layer]
\label{def:observer}
A cognitive observer adds a representation map $\Phi \colon W \to Y$ from the
world's task structure $W$ to the consequence space on which the evaluation
object is defined, and a search policy $\pi$ over represented states. The
observer is the quadruple $(\Phi, G, B, K)$ together with $\pi$. The world
retains $W$ and the action set.
\end{definition}

Definition~\ref{def:observer} is where OT enters. $\Phi$ is a read operator, and
OT's finding is that its content is a property of the pair and not of the world.
The cognitive reading is that two agents facing the same task can face different
operational problems, and that this happens before anything that could be called
a preference.

\section{The vacuity problem, stated as a proposition}

The programme's own statement of limits names flexibility as the principal
threat. We state it more sharply, because a threat that is described in prose
can be acknowledged indefinitely without being retired.

\begin{proposition}[Vacuity without modular identification]
\label{prop:vacuity}
Let a finite choice record consist of menus $M_1,\dots,M_n$ drawn from a finite
alternative set and a selection from each. If the observer's four components
$(\Phi, G, B, K)$ are all fitted to that record, then for any record there
exists an observer reproducing it exactly. The account therefore excludes no
finite choice record, and no observation of choice alone can disconfirm it.
\end{proposition}

\noindent\emph{Status: proved, elementary, and not new.} Admissibility alone
suffices. Taking $K$ menu-dependent and equal to the singleton actually chosen
reproduces any record whatever, for any $\Phi$, $G$ and $B$.

The proposition belongs to a known family and is restated here only because the
components are ours. Kalai, Rubinstein and Spiegler show that rationalization by
multiple rationales admits every choice function, so the model measures degrees
of rationality rather than rationality \cite{krs2002}. Under endogenous menu
assignment every dataset of menu frequencies and conditional choice
probabilities is rationalizable by some random-utility distribution, and the
classical axioms lose their content entirely. The contrast with Afriat is the
instructive part \cite{afriat1967}. On budget data, rationalizability by a
concave monotone utility has real content, because GARP can fail. What removes
the content in our setting is menu-dependent admissibility, which is exactly the
component GET keeps separate from cost because it can violate the weak axiom.
The separation that makes $K$ theoretically interesting is the same separation
that makes the account untestable when $K$ is free.

The proposition's interest is therefore not mathematical. It is that vacuity
arrives through the weakest of the four components, so no argument about the
plausibility of a metric or the realism of a budget repairs it.

Proposition~\ref{prop:vacuity} has the shape of a defect this programme has met
before in another arm. A conformance property evaluated over an empty dependency
graph passes perfectly and certifies nothing, and the repair there was to
require that the structure the property quantifies over be non-vacuous before
the property may be reported. The repair here is the same in form. The
components of the observer must be identified on evidence that the prediction
does not use.

\begin{definition}[Modular identification]
\label{def:modular}
An OLGC account of a choice record $R$ is \emph{modularly identified} when each
of $\Phi$, $G$, $B$ and $\pi$ is estimated from evidence disjoint from $R$, and
$K$ is either fixed in advance or derived from a stated constraint that is not a
function of the selections in $R$.
\end{definition}

\begin{proposition}[Content under modular identification]
\label{prop:content}
Under Definition~\ref{def:modular} with $\Phi$, $G$ and $B$ fixed before $R$ is
opened, the induced order on any menu is determined, and the account therefore
excludes every record that differs from the induced one by more than the
semiorder's own indifference band.
\end{proposition}

\noindent\emph{Status: proved, given GET's Theorem 2.} The band is $\varepsilon_B$
and is fixed by the budget rather than fitted, which is the property that makes
the exclusion non-empty. The proposition is the reason the programme's
experimental designs insist on separate calibration blocks. It converts a
methodological preference into a condition for the theory to have content.

A second repair exists and does not require external identification. Caplin and
Dean give a revealed-preference test for optimal acquisition of costly
information whose no-improving-attention-cycles condition is necessary and
sufficient for a finite dataset to be consistent with the model, and they
recover information costs from choice data alone \cite{caplindean2015}. Related
work identifies consideration sets from choice in the same spirit
\cite{masatlioglu2012,manzini2007}. That line and this one answer the same
threat differently. Theirs constrains the model class until choice data
discriminate within it. Ours holds the model class and constrains the evidence.
Which is preferable is an empirical question about how much external
identification costs, and it is not settled here. A budget recovered by a
cycle condition and a budget manipulated experimentally are different objects,
and P1 below is a test of the second.

\begin{open}
\label{open:partial}
Proposition~\ref{prop:content} assumes full prior identification. The realistic
case is partial, where one component is estimated jointly with the prediction.
The degrees of freedom retired by identifying $\Phi$ alone, or $B$ alone, are
not characterized, and no bound is known on how much a single free component can
absorb. Until such a bound exists, an account that identifies three of four
components is not known to be more falsifiable than one that identifies none.
\end{open}

\subsection*{Identification can also be too good}

Open~\ref{open:partial} is one end of a problem with two ends. Under-identified,
a free component absorbs the prediction and nothing is excluded. The other end
appeared in the Phase I pilot and was not anticipated when
Proposition~\ref{prop:content} was written.

The Phase I design identifies $G$ and $t$ from a calibration block of ordinal
comparisons, then admits a pair into the trading class precisely when the fitted
metric says the pair's rank-$k$ order disagrees with its full-budget order, and
into the within-subspace class precisely when it says the order is preserved.
Reversal is predicted for the first class and not for the second. The class
label is therefore the prediction, and the consequence is not approximate. Let
$A$ be the event that the evaluator's order agrees with the fitted metric at
full budget and $B$ that it agrees at the reduced budget. A trading pair is one
the metric says reverses, so the evaluator reverses it exactly when $A$ and $B$
take the same value; a within-subspace pair is one the metric says is preserved,
so the evaluator reverses it exactly when they differ. Hence
\[
R_T = \Pr(A = B), \qquad R_W = 1 - \Pr(A = B), \qquad
R_T - R_W = \Pr(A=B \mid T) + \Pr(A=B \mid W) - 1 .
\]
The contrast the design grades is a function of two agreement rates and of
nothing else. Measured on the pinned evaluator across six cells, the marginals
reproduce every observed reversal rate to a maximum error of $0.0074$, against
an instrument reproducibility floor of $0.0070$ established by rerunning
identical pairs. The correlation between $A$ and $B$ was measured rather than
assumed and is zero, at $-0.053$ and $-0.068$ in the two cells where both
marginals have variance, so the independent form
$R_T = p_A p_B + (1-p_A)(1-p_B)$ is exact. No part of the separation is left for
the budget manipulation to explain, and the pass bars are the same two rates
restated.

This is not an artifact of one design. Any test of the form \emph{identify the
components, then predict a choice from them} inherits the difficulty, because
the prediction is a function of the identified components and its
informativeness falls as the identification error falls toward zero. Content is
therefore not monotone in identification accuracy. It is low at both ends and
the maximum is somewhere in between, and we can locate neither end.

The same reread found a second condition, on $B$ rather than on $G$, which
Definition~\ref{def:modular} does not state and should. A rank budget is a
choice of top-$k$ eigenspace of the workload moment
$M = \Eb[\,G^{1/2}(C-t)(C-t)^{\!\top} G^{1/2}\,]$, and that eigenspace is
identified only if $M$ has a spectral gap at $k$. The Phase I design drew
consequences uniformly on a cube and placed the ideal at the cube's centroid,
which makes $M$ isotropic analytically, at $(\mathrm{hi}-\mathrm{lo})^2/12$
times the identity. There is then no gap and no preferred eigenspace, and the
retained subspace is fixed by which finitely many calibration pairs were drawn.
On an evaluator that is exactly Euclidean, with no response error whatever, two
independent calibration draws select retained directions a median of $56$
degrees apart, against the $60$ degrees expected between two uniformly random
directions in three-space. Under that design a budget manipulation is a
projection onto an arbitrary plane, and a confirmation of P2 would report a
property of orthogonal projection. Identifying $\Phi$ and $G$ on disjoint
evidence does not help, because the defect is in the workload distribution the
moment is taken over.

\begin{definition}[Budget identification]
\label{def:gap}
A rank budget $B = k$ is identified for an evaluation object only if the
workload moment $M$ has a spectral gap at $k$, and an experiment that
manipulates $B$ must exhibit that gap against the sampling distribution of $M$
at the realized sample size rather than assume it.
\end{definition}

\noindent The condition is easy to satisfy and easy to omit, and omitting it
costs the whole experiment. It also shows that the anti-vacuity check the
programme inherited from its sibling arm is not sufficient here. That check
requires a non-trivial share of the trace to be discarded, and isotropy
discards the largest such share while distinguishing nothing.

\begin{open}
\label{open:band}
Proposition~\ref{prop:content} treats prior identification as a good whose
quantity is to be maximized. The pilot shows it is not. No bound is known on how
accurate an identification may be before the prediction it licenses becomes a
consequence of the fit, and no measure separates the residual content from the
entailed part. Until one exists, a design of this shape should register
identification accuracy as an interval and not as a floor, and should expect the
informative work to be done by the identification step rather than by the
prediction it feeds.
\end{open}

The consolation is that the identification step is where the failures actually
occurred. Of three candidate evaluators put through the Phase I calibration,
two could not be calibrated at all. One ordered by presentation position at a
first-position rate of $0.8175$, with an agreement rate across the two orders
that matched what pure position preference predicts and nothing that a metric
would explain. A second recovered distances at $r^2 = 0.766$ and still ordered
by position at $0.700$. Only the third compared by content. An evaluator can
fail to be a quadratic evaluation object, and most of the ones we tried did, so
the step that carries the content is the step that can fail.

\section{Manipulating a budget, and three ways it failed}
\label{sec:manip}

Definition~\ref{def:gap} says when a rank budget is identified. It does not say
how to change one, and the first phase of this programme spent three designs
discovering that the difference matters. All three ran against the same pinned
language-model evaluator. None produced a resolution effect, and they failed for
three unrelated reasons rather than one. An earlier version of this section
grouped the second and third as a common failure mode; the records refute that,
since the second condition's unusable verdicts all parsed and disagreed while the
third's had mostly not parsed at all. They are reported because they are more
informative than the designs would have been had they worked, and because the
first two failure modes are available to any experiment of this shape.

\subsection*{A projected stimulus is not a reduced budget}

The first design rendered each option's rank-$k$ projection and showed the
evaluator those numbers. This looks like a budget manipulation and is not one.
Theorem 5 of GET concerns \emph{fixed} actions evaluated under \emph{different}
budgets; projecting the rendering holds the budget fixed and changes the actions.
The evaluator was never resolution-limited. It read what it was handed, which is
why its order agreed with the rank-$k$ order in four of six cells at a rate of
exactly $1.000$.

The consequence is the identity in \S3. With the class label read off the same
fitted metric that generates the prediction, and the stimulus pre-projected so
that agreement at the reduced budget is automatic, the graded contrast reduces to
two agreement rates. Across six cells the marginals reproduced every observed
reversal rate to a maximum error of $0.0074$, against an instrument
reproducibility floor of $0.0070$ established by rerunning identical pairs. The
design could not have failed and therefore could not have succeeded.

\subsection*{Removing the evaluator's deliberation removes the instrument}

The second design held the stimulus fixed at full precision and varied the
evaluator, disabling the chain of thought so that comparisons ran at zero
reasoning tokens rather than about $350$. That is a genuine resource
manipulation. It is also not a coarsening.

Under the manipulation the evaluator stopped comparing by content and began
choosing by presentation position. Pairs are graded only when the verdict
survives exchanging the two options, and between $1.6$ and $10.9$ percent of
pairs survived, against $98$ percent with deliberation intact. Writing $p$ for
the probability of choosing by position, survival is about $1-p$, which places
the condition between $89$ and $98$ percent position-driven. The programme had
already recorded the same pattern in two other models, at first-position rates
of $0.8175$ and $0.700$. This condition produced no unparsed answers at all, so
its unusable verdicts are genuine order-dependence and not a generation limit,
which is what distinguishes it from the third design below. Deliberation is what makes this class of evaluator
compare by content at all, so removing it removes the measuring instrument rather
than narrowing the thing measured.

\subsection*{Load degrades reliability before it degrades resolution}

The third design kept deliberation and added distractor attributes, drawn per
pair and identical across the two options of that pair. Their contribution to
$d(a)^2-d(b)^2$ is then exactly zero, so the correct answer is invariant by
cancellation at any load, while the evaluator must still read and weigh the
added coordinates. At sixteen distractors the shared term adds about $34{,}000$
to both distances and the difference between them is unchanged.

This manipulation kept the evaluator comparing by content. First-position rates
stayed flat across the ladder, and reasoning rose from $398$ to $918$ tokens, so
the load landed. Scoring every pair at every load, with an ungradeable verdict
counted as uninformative rather than dropped, the trading-pair score was flat,
and its intervals never approached one half from above.

That null is weaker than it first appeared, for two reasons found by a later
audit of the record rather than by the design. Reasoning at the highest load
averaged $918$ tokens against a generation limit of $1024$, and the scoring
treated an answer that never parsed identically to a verdict that parsed in both
presentation orders and disagreed. The two counts track each other almost
exactly. Bounding the difference, no cell at the highest load contains more than
four genuine disagreements in sixty-four. What looked like a decline in the
evaluator's reliability under load was the generation limit, and the apparent
descent of the control class toward the trading class was the same artifact.
Second, the scoring convention was selected after the fact from three that had
been computed: over the interval the trading score moves by $+0.000$ counting an
ungradeable verdict as uninformative, by $-0.094$ counting it as an error, and by
$+0.073$ conditioning on gradeable pairs alone. Three conventions, three signs.

The deeper gap is that nothing in this design establishes that the load reduced
the evaluator's effective rank. The distractor coordinates are visibly
uninformative, and an evaluator that sheds uninformative dimensions first is not
exhibiting the budget the theory describes, which discards directions that are
informative but low in variance. No spectrum was recalibrated under load. A null
result from a manipulation that was never shown to move the quantity it is named
after is not evidence about budgets.

\subsection*{The estimand has to be chosen before the manipulation is run}

One methodological point generalizes beyond this evaluator, and it inverted a
result before it was caught. An analysis that conditions on a filter which the
manipulation itself moves will report the effect backwards. Scoring accuracy over
the pairs that survived the swap check, rather than over all pairs, made trading
accuracy appear to \emph{rise} with load, from $0.857$ to $0.930$. It rose because
the pairs lost to inconsistency were disproportionately ones the evaluator had
been getting wrong: under random attrition nine errors among sixty-three graded
pairs would leave about six among forty-three survivors, and three were observed,
with a hypergeometric probability of $0.023$. On a fixed denominator the same
data show no trend at all. Any design in which a manipulation can affect which
observations are gradeable needs its estimand fixed in advance over the full set.

\begin{open}
\label{open:manipulable}
None of the three manipulations produced a resolution effect, and two failed for
reasons that would recur in any similar design. Whether a language-model
evaluator has a rank budget that can be manipulated at all is therefore open. The
alternatives are not distinguished by this evidence: that the budget exists and
no intervention tried was the right one, that the resource such an evaluator
spends is elastic rather than fixed so that a heavier task is met with more
computation rather than a coarser geometry, or that the rank-budget abstraction
does not describe this substrate. The evidence is weaker than three failed
attempts suggests, because no attempt included a manipulation check. Nothing
measured whether any intervention reduced effective rank, so a design that
recalibrates the metric under the manipulation and reports the spectrum is a
precondition for the next attempt rather than a refinement of it. Until a manipulation degrades resolution while
leaving swap consistency intact, P2 is untested here rather than unsupported.
\end{open}

\section{What is proved, what is measured, and what is posited}

The programme's statements do not carry a uniform status, and a foundational
paper that presents them at one strength misrepresents all of them.

\begin{table}[!ht]
\caption{Status of the programme's statements. The evaluation layer is proved,
and measured on artificial evaluators except where a row says otherwise. Every
statement specific to \emph{cognition} sits in the lower block and carries no
sealed measurement.}
\label{tab:status}
\centering\footnotesize
\begin{tabular}{@{}p{0.30\textwidth} p{0.14\textwidth} p{0.44\textwidth}@{}}
\toprule
\textbf{Statement} & \textbf{Status} & \textbf{Basis} \\
\midrule
Rank budget coarsens distinctions & proved & GET Theorem 1, machine-checked \\
Preference is a semiorder with threshold the budget & proved & GET Theorem 2, machine-checked \\
Budget change reverses preference between fixed actions & proved (existence) & GET Theorem 5, machine-checked \\
Threshold tracks budget in an evaluator & demonstrated & GET gate G3, sealed. A language-model scorer's indifference threshold is 0.87 to 0.88 of the rendering step, and every pair separated by 20 is ordered alike in every cell \\
Hull law in a population & demonstrated as a tendency & GET gate G2, sealed, ANES 1972. Half of 489 testable respondents violate at least once against a shuffle null of 66 percent. It fails as a deterministic statement \\
An evaluator can fail to be a quadratic evaluation object & measured, not sealed & Phase I instrument search \cite{c1}. Two of three candidate language-model evaluators could not be calibrated, ordering by presentation position at 0.8175 and 0.700. The third calibrated and held out at 0.985, which is a maximum selected over the three and should be read as one \\
A rank budget in a language-model evaluator is manipulable & open, three attempts, no effect & \S\ref{sec:manip} and Open~\ref{open:manipulable}. One design did not manipulate a budget, one destroyed the instrument, one preserved it and produced no resolution effect \\
Representation is prior to preference & posited for cognition & measured in other domains by OT, not in a cognitive observer \\
Human cognitive budgets exist and are manipulable & posited & no sealed measurement \\
Bounded search is the right dynamical abstraction & posited & alternatives in \S6 are not excluded \\
Whitney stratification of cognitive state space & posited & no measurement distinguishes it from a steep smooth family \\
Cross-architecture universality & posited & no collapse has been attempted \\
\bottomrule
\end{tabular}
\end{table}

The four rows at the bottom are the content of the synthesis. They are also the
rows with no sealed measurement behind them. That is the position of the
programme at draft 0.1 and it is stated here rather than in a closing section.

\section{Predictions, as consequences of the identification conditions}

The programme states five predictions. Each is restated here as a claim that
becomes testable only under Definition~\ref{def:modular}, with the component
whose independent identification it requires named beside it.

\textbf{P1, the threshold-budget law.} Holding $G$ and $t$ fixed and
manipulating $B$, the just-noticeable comparison threshold tracks the budget.
Well separated pairs keep their order. Requires independent identification of
$G$ and $t$. A systematic reordering of well separated pairs under a pure
resolution manipulation contradicts the model.

\textbf{P2, direction-specific rank reversal.} Reversals concentrate on pairs
trading a retained direction against a newly unresolved one. Pairs whose
difference lies in the retained subspace keep their order. Requires independent
identification of $\Phi$, since the retained subspace is a property of $\Phi$
and $B$ together. This is the prediction that separates the account from a claim
that load adds noise. Definition~\ref{def:gap} and Open~\ref{open:band} bound what a
confirmation of it is worth. The retained subspace must be identified before it
can be manipulated, and where the class assignment is made by the same fitted
metric that predicts the reversal the contrast is an identity rather than a
test.

\textbf{P3, cross-measure boundary alignment.} At a genuine regime transition,
several observables reorganize near the same control value. Requires a boundary
located from one observable and registered before the others are examined.
Without the registration the prediction is unfalsifiable in practice, because a
boundary can be located in any observable after the fact.

\textbf{P4, representation-first generalization.} An account whose $\Phi$ is
identified on one task family and whose $G$ is identified on another predicts a
third without refitting. This is Definition~\ref{def:modular} stated as an
experiment rather than as a condition.

\textbf{P5, cross-architecture universality.} Normalized quantities such as a
margin-to-budget ratio collapse across substrates. Requires the normalization to
be fixed before the second substrate is measured. Failure weakens the claim to a
fundamental theory without refuting the evaluation layer.

\section{Falsifiers, as conditions on the record}

The programme lists seven falsifiers. Stated as conditions on a record rather
than as judgments, they are checkable by a reader who did not run the
experiment.

\begin{enumerate}[label=F\arabic*.]
\item A registered budget manipulation leaves the distinction threshold and the
predicted reversal classes unchanged.
\item A representation probe registered in advance shows no change where the
account attributes a flip to representation.
\item Choices reverse among alternatives whose differences lie in a subspace
shown to be retained, by a probe registered before the choices are collected.
\item A scalar model with fewer effective degrees of freedom predicts the same
held-out reversals and process measures at least as well.
\item Registered regime boundaries fail to align across independent observables.
\item Cross-domain metric transfer disappears when researcher-coded encodings
are replaced by independently measured ones.
\item No reusable relation between representational limit and behavioral
transition appears across substrates.
\end{enumerate}

F6 is the one the existing domain exposition already fails in its current form,
because its encodings are researcher-coded. That is recorded rather than
defended. Phase II exists to answer it.

\section{Alternatives that are not excluded}

A foundational paper that lists no live alternative is not reporting the state
of the question.

Geometry may be descriptive rather than mechanistic. A successful geometric
account of endpoints does not establish that the process follows geometric
paths, and the domain exposition is explicit that its basis is a modeling
vocabulary. Distinguishing the two requires process evidence, which Phase III
exists to collect.

Stratification may be too strong. A steep smooth family approximates most
categorical transitions. The stratified model earns its place only by predicting
a regularity at the boundary that the smooth family misses, and no such
prediction has been registered.

Search may be the wrong abstraction. Sampling, diffusion, predictive coding and
direct policy execution are all live. The account survives only if bounded
search is read at a level of abstraction where some counterfactual or trajectory
structure is measurable, and if nothing of that kind is measurable the dynamic
layer should be reformulated rather than defended.

\section{Related work}
\label{sec:related}

The components of this account are borrowed and the borrowing is stated here
rather than in a footnote.

Finite resolution producing intransitive indifference is Luce's semiorder
\cite{luce1956}, and the behavioral fact it models is older than the
formalism \cite{tversky1969}. Bounded resources producing satisficing rather
than optimization is Simon's \cite{simon1955}. The result that finite data
underdetermine the model is Afriat's on budget data \cite{afriat1967} and
Kalai, Rubinstein and Spiegler's on choice functions \cite{krs2002}. Recovering
an attention or information cost from choice alone is Caplin and Dean's
\cite{caplindean2015}, with consideration-set identification in
\cite{masatlioglu2012,manzini2007}.

What this paper adds is narrow. It states the vacuity of the specific
four-component observer as a proposition rather than as a limitations
paragraph, gives modular identification as a condition on evidence rather than
as methodological advice, and organizes a phase plan by which degree of freedom
each phase retires. We do not offer a better theory of choice. We offer the
condition under which this one can be wrong.

\section{Programme}

Each phase retires one degree of freedom named in
Proposition~\ref{prop:vacuity}. That is the organizing principle, and a phase
that retires none is not a phase.

\begin{table}[!ht]
\caption{The phase plan, organized by the degree of freedom each phase retires
from Proposition~\ref{prop:vacuity}. A phase that retires none is not a phase.
Only Phase I has a fixed design at this draft.}
\label{tab:phases}
\centering\footnotesize
\begin{tabular}{@{}p{0.06\textwidth} p{0.40\textwidth} p{0.44\textwidth}@{}}
\toprule
& \textbf{Objective} & \textbf{Degree of freedom retired} \\
\midrule
I & A causal observer-limited reversal under registered budget manipulation & $B$ becomes manipulated rather than fitted \\
II & Representation recovered from independent tasks or learned encoders & $\Phi$ stops being researcher-coded, answering F6 \\
III & Trajectories from process measures & $\pi$ becomes measured rather than assumed \\
IV & A registered boundary with a prediction the smooth family misses & stratification becomes an empirical claim or is dropped \\
V & A collapse of normalized observables across substrates & the word fundamental acquires a referent or is withdrawn \\
VI & A safety-relevant boundary predicted and prevented without training on the failure & the account becomes an instrument \\
\bottomrule
\end{tabular}
\end{table}

Phase I is deliberately low-dimensional and is the only phase whose design is
fixed at this draft. Its registration fixes the consequence vectors, the
candidate retained subspaces, and the predicted ordering at each budget before
any data are opened \cite{c1}. The prediction it tests is P2, because P2 is the
one that a noise account does not also make.

Phase I is designed but not sealed, and the distinction is worth keeping. Its
evaluator is a hosted language model pinned by alias and creation timestamp,
with a preflight that refuses to run if either has moved, because the gateway
publishes no revision and an alias can be repointed without the identifier
changing. That converts an unverifiable provenance claim into a check that fails
loudly rather than into a sentence in a methods section. Its pass tolerances
were fixed by a pilot on a seed disjoint from the run's, at a reversal margin of
$0.484$ and an ambiguity tolerance of $0.05$, with the rendering ceiling not
binding. The pilot is spent. It grades no claim, its numbers are not revised,
and the run seed is drawn only after the registration is sealed. What it has
established is that the design is executable on a real evaluator and that its
tolerances were not chosen with the answer in view.

\section{What this paper has not done}

No gate in this paper has been sealed, and Phase I is not sealable as
registered. A cold reread of its registration returned three defects that change
what the experiment measures, of which Definition~\ref{def:gap} and the identity
in Section~3 are two, together with five registered bars that cannot fire and an
unregistered line in the sampler that decides a registered outcome \cite{c1}.
The three manipulations in \S\ref{sec:manip} are design work and
grade nothing either, and their common failure, that a resource intervention
reaches the instrument before it reaches the geometry, is a fact about this
substrate and this instrument rather than a result about budgets.
Open~\ref{open:manipulable} is unresolved and bounds P2 in this programme: the
prediction is untested here, not unsupported. The Phase I pilot and the
instrument search are design work and grade nothing, and the one row of
Table~\ref{tab:status} they support is marked unsealed. Open~\ref{open:partial}
is unresolved and bounds every claim about partial identification.
Open~\ref{open:band} is unresolved and bounds what a Phase I pass would be
worth. No measurement of a human budget exists in this programme, and nothing
here licenses carrying a language-model result across to one. The domain
exposition's measurements are retained at the class its own admission filter
assigned them and are not promoted by the appearance of this paper.

\section*{Acknowledgment of AI use}

Drafting and structural editing were performed with the assistance of a large
language model under the author's direction. The claims, the status assignments,
and the decision to state vacuity as a proposition rather than as a limitation
are the author's.

\begin{thebibliography}{9}
\footnotesize
\bibitem{get} A.~H. Bond, \emph{Geometric Evaluation Theory}, foundational paper
draft 0.2, San Jos\'e State University, 2026. Claim ledger at
\texttt{claims/LEDGER.md}. % verify draft number at submission
\bibitem{ot} A.~H. Bond, \emph{Observation Theory}, campaign record,
\texttt{ahb-sjsu/observation-theory-campaigns}, 2026.
\bibitem{gc} A.~H. Bond, \emph{Geometric Cognition, The Mathematical Structure
of Human and Artificial Thought}, draft, 2026.
\bibitem{c1} A.~H. Bond, \emph{C1, Direction-Specific Rank Reversal in a
Language-Model Evaluator}, Phase I registration and stage records,
\texttt{experiments/C1/}, 2026. Unsealed at the time of writing.
\bibitem{thesis} A.~H. Bond, \emph{Observer-Limited Geometric Cognition, A
Candidate Fundamental Theory of Cognition from Representation, Evaluation, and
Bounded Search}, research thesis, September 2026.
\bibitem{afriat1967} S.~N. Afriat, The construction of utility functions from
finite demand data, \emph{International Economic Review}, 8(1), 67--77, 1967.

\bibitem{krs2002} G.~Kalai, A.~Rubinstein and R.~Spiegler, Rationalizing choice
functions by multiple rationales, \emph{Econometrica}, 70(6), 2481--2488, 2002.
% verify issue and pages on final

\bibitem{caplindean2015} A.~Caplin and M.~Dean, Revealed preference, rational
inattention, and costly information acquisition, \emph{American Economic
Review}, 105(7), 2015.
% verify pages on final; do not copy Masatlioglu's range, they are different papers

\bibitem{masatlioglu2012} Y.~Masatlioglu, D.~Nakajima and E.~Y. Ozbay, Revealed
attention, \emph{American Economic Review}, 102(5), 2012.
% verify pages on final

\bibitem{manzini2007} P.~Manzini and M.~Mariotti, Sequentially rationalizable
choice, \emph{American Economic Review}, 97(5), 1824--1849, 2007.

\bibitem{luce1956} R.~D. Luce, Semiorders and a theory of utility discrimination,
\emph{Econometrica}, 24(2), 178--191, 1956.
\bibitem{tversky1969} A.~Tversky, Intransitivity of preferences,
\emph{Psychological Review}, 76(1), 31--48, 1969.
\bibitem{simon1955} H.~A. Simon, A behavioral model of rational choice,
\emph{Quarterly Journal of Economics}, 69(1), 99--118, 1955.
\end{thebibliography}

\end{document}
